\documentclass[english]{lapesd-thesis}

% Outros \usepackage{}
\usepackage{capt-of}
\usepackage{pgfgantt}
% \usepackage{rotating}   % sidewaystable
% \usepackage{ragged2e}   % RaggedRight nas colunas
% \newcolumntype{L}[1]{>{\RaggedRight\arraybackslash}p{#1}}
\usepackage{longtable}
\usepackage{booktabs}
\usepackage{array}


%%%%%%%%%%%%%%%%%%%%%%%%%%%%%%%%%%%%%%%%%%%%%%%%%%%%%%%%%%%%%%%%%%%%
%%% Configurações da classe (dados do trabalho)                  %%%
%%%%%%%%%%%%%%%%%%%%%%%%%%%%%%%%%%%%%%%%%%%%%%%%%%%%%%%%%%%%%%%%%%%%

% Informações para capa e folha de rosto/certificacao
\titulo{An Orchestration Architecture for Deployment and Management of Digital Twins in Industrial Environments}
\autor{André Roberto Alves de Oliveira}
% \data{1 de agosto de 2019} % ou \today
\data{\today} % ou \today
% \tese % ou \dissertacao
\dissertacao
\preambulo{Trabalho de qualificação submetido ao Programa de
Pós-Graduação em Ciência da Computação como requisito parcial
para a obtenção do título de Mestre em Ciência da Computação.}
\titulode{Mestre em Ciência da Computação}
\orientador{Prof. Frank Augusto Siqueira, Dr.}
% \coorientador{Prof. Lars Thørväld, Dr.}

%%% Atenção! No caso de TCC, além de usar \tcc, outros comandos devem ser 
%%% fornecidos:
% \tcc
% \departamento{Departamento de Informática e Estatística}
% \curso{Ciência da Computação}
% \titulode{Bacharel em Ciência da Computação}
% %% Para TCCs, orientadores e coorientadores podem ser externos, logo a
% %% BU exige que sua afiliação seja explicitada. Por padrão, assume-se
% %% UFSC. Você pode alterar a afiliação com os comandos abaixo:
% \afiliacaoorientador{Universidade Federal de Santa Catarina}
% \afiliacaocoorientador{Universidade Federal da Terra de Ninguém}

% Membros da banca e coordenador
% As regras da BU agora exigem que Dr. apareça depois do nome
% Dica: para gerar Profᵃ. use Prof\textsuperscript{a}.
% Dica 2: para feminino use \orientadora e \coorientadora
\membrobanca{Prof. Rodrigo Pereira, Dr.}{Universidade Federal de Santa Catarina}
\membrobanca{Prof. Patrícia Plentz, Dr.}{Universidade Federal de Santa Catarina}
\membrobanca{Prof. Eduardo Camilo Inácio, Dr.}{Universidade Federal de Santa Catarina}
% \coordenador{Prof. Charles Palmer, Dr.}


% Ativa indíce remissivo. Precisa estar aqui, não funciona no .cls nem no
% \BeforeBeginDocument{}
\makeindex
% Carrega definições dos acrônimos e do glossário. Isso precisa ser feito antes
% do \begin{document}

%%%%%%%%%%%%%%%%%%%%%%%%%%%%%%%%%%%%%%%%%%%%%%%%%%%%%%%%%%%%%%%%%%%%
%%% Lista de acronimos                                           %%%
%%%%%%%%%%%%%%%%%%%%%%%%%%%%%%%%%%%%%%%%%%%%%%%%%%%%%%%%%%%%%%%%%%%%
%%% Importante:                                      
%%% - A lista PRECISA SER MANTIDA ORDENADA
%%%%%%%%%%%%%%%%%%%%%%%%%%%%%%%%%%%%%%%%%%%%%%%%%%%%%%%%%%%%%%%%%%%%

\tnewacronym{API}{Application Programming Interface}
            % \API e \APIs nunca espandem, mas ele aparece na lista de siglas
\xnewacronym{DHT}{Distributed Hash Table}
            % Se \DHTs for o primeiro uso, Tabe ganha um s no final
\xnewacronym[][longplural={Square Matrices}]{SQ}{Square Matrix}
            % trata o plural usado em \SQs
            % note que como queremos passar o 2o arg opcional, devemos 
            % passar o primeiro (podemos deixar em branco)
\xnewacronym[WTC]{W3C}{World Wide Web Consortium}
            % gera comando \WTC que expande para "World ... (W3C)"

\xnewacronym{AAS}{Asset Administration Shell}
\xnewacronym{ADT}{Atomic Digital Twin}
% \xnewacronym[][longplural={Application Programming Interfaces}]{API}{Application Programming Interface}
\xnewacronym{BoD}{Bill of Disposal}
\xnewacronym{BoM}{Bill of Materials}
\xnewacronym{BoP}{Bill of Process}
\xnewacronym{BoS}{Bill of Services}
\xnewacronym{CN}{Connections in Digital Twin}
\xnewacronym{CDT}{Composite Digital Twin}
\xnewacronym{CPS}{Cyber-Physical System}
\xnewacronym{DaaS}{Devices-as-a-Service}
\xnewacronym{DB}{Database}
\xnewacronym{DSR}{Design Science Research}
\xnewacronym{DD}{Digital Twin Data}
\xnewacronym{DT}{Digital Twin}
\xnewacronym{DTaaP}{Digital Twin as a Platform}
\xnewacronym{DTaaS}{Digital Twin as a Service}
\xnewacronym{DTC}{Digital Twin Client}
\xnewacronym{DTI}{Digital Twin Instance}
\xnewacronym{DTM}{Digital Twin Master}
\xnewacronym{DTN}{Digital Twin Network}
\xnewacronym{DTP}{Digital Twin Prototype}
\xnewacronym{DTE}{Digital Twin Environment}
\xnewacronym{DTPod}{Digital Twin Pod}
\xnewacronym{EC}{European Commission}
\xnewacronym{EDT}{Edge Digital Twin}
\xnewacronym[][longplural={European Interoperability Frameworks}]{EIF}{European Interoperability Framework}
\xnewacronym[][longplural={Geometric Dimensioning and Tolerancings}]{GDT}{Geometric Dimensioning and Tolerancing}
\xnewacronym[GDNT][long={Geometric Dimensioning and Tolerancing}, short={GD\&T}, sort={GDT}]{GDNT}{Geometric Dimensioning and Tolerancing}
\xnewacronym{IoT}{Internet of Things}
\xnewacronym{IoTw}{Internet of Twins}
\xnewacronym{KPI}{Key Performance Indicator}
\xnewacronym{LCIM}{Levels of Conceptual Interoperability Model}
\xnewacronym{MAS}{Multi-Agent System}
\xnewacronym{MQTT}{Message Queuing Telemetry Transport}
\xnewacronym{NIST}{National Institute of Standards and Technology}
\xnewacronym[OPCUA][long={OPC-UA}]{OPC-UA}{OPC Unified Architecture}
\xnewacronym{PE}{Physical Entity}
\xnewacronym{PLM}{Product Lifecycle Management}
\xnewacronym[][longplural={REpresentational State Transfer APIs}]{REST}{REpresentational State Transfer}
\xnewacronym{Ss}{ervices in Digital Twin}
\xnewacronym[SthreeStorage][long={Simple Storage Service}, short={S3}, sort={S3}]{SthreeStorage}{Simple Storage Service}
\xnewacronym{SoS}{System of Systems}
\xnewacronym{VM}{Virtual Model}
\xnewacronym{XaaS}{Everything as a Service}


%%% Local Variables:
%%% mode: latex
%%% TeX-master: "main"
%%% End:

\input{glossary.tex}

\begin{document}

%%%%%%%%%%%%%%%%%%%%%%%%%%%%%%%%%%%%%%%%%%%%%%%%%%%%%%%%%%%%%%%%%%%%
%%% Conteúdo                                                     %%%
%%%%%%%%%%%%%%%%%%%%%%%%%%%%%%%%%%%%%%%%%%%%%%%%%%%%%%%%%%%%%%%%%%%%

\pretextual%
% Assume-se que \pretextual já foi feito


\imprimircapa%
\imprimirfolhaderosto*
% Atenção! esse \protect é importante
\protect\incluirfichacatalografica{ficha.pdf}
\imprimirfolhadecertificacao


\begin{dedicatoria}
  Este trabalho é dedicado à wikipedia e ao stackoverflow. Essa frase está aqui apenas para testar o alinhamento e as margens
\end{dedicatoria}


\begin{agradecimentos}
  O presente trabalho foi realizado com apoio da Coordenação de Aperfeiçoamento de Pessoal de Nível Superior -- Brasil (CAPES) -- Código de Financiamento 001.
\end{agradecimentos}


\begin{epigrafe}
  For a number of years I have been familiar with the observation that the quality of programmers is a decreasing function of the density of go to statements in the programs they produce \\
  \cite{dijkstra1968}
\end{epigrafe}


\begin{resumo}[Resumo]
  Aqui deve ser inserido um resumo de 150 a 500 palavras (limitação de tamanho dada pela BU). A linguagem deve ser português e a hifenização já foi alterada. O resumo em português deve preceder o resumo em inglês, mesmo que o trabalho seja escrito em inglês. A BU também diz que deve ser usada a voz ativa e o discurso deve ser na 3ª pessoa. A estrutura do resumo pode ser similar a estrutura usada em artigos: Contexto -- Problema -- Estado da arte -- Solução proposta  -- Resultados.

  % Atenção! a BU exige separação através de ponto (.). Ela recomanda de 3 a 5 keywords
  \vspace{\baselineskip} 
  \textbf{Palavras-chave:} Palavra-chave. Ponto como separador. Bla.
\end{resumo}


\begin{resumo}[Resumo Estendido]
%%%%%%%%%%%%%%%%%%%%%%%%%%%%%%%%%%%%%%%%%%%%%%%%%%%%%%%%%%%%%%%%%%%%%%
% Atenção: normas e templates contraditórios!!!                    %%%
%%%%%%%%%%%%%%%%%%%%%%%%%%%%%%%%%%%%%%%%%%%%%%%%%%%%%%%%%%%%%%%%%%%%%%
% - Modelo da BU: https://repositorio.ufsc.br/handle/123456789/197458
% - A BU exige no **mínimo** 2 páginas e no **máximo** 5
% - Regimento do PPGCC, Art 40 Entende-se  por  resumo  estendido  um  documento  que  contenha  as  informações  mais  relevantes  de  cada  capítulo  da  tese  ou  da  dissertação.
% O mais seguro é ignorar o regimento e seguir a BU.
    % Atenção! A BU diz que o resumo **deve** conter as seções abaixo!
  \section*{Introdução} % Deve ser  subsection*, devido a formatação usada no modelo
  A hifenização é alterada para \texttt{brazil}, mesmo para documentos em inglês. Descrever brevemente esses itens exigidos pela BU. Como a RN 95/CUn/2017 é mais recente e impõe outras regras a revelia de regimentos e regulamentos, é mais sábio obedecê-la. Lembre que esse resumo estendido deve term entre 2 e 5 páginas.
  
  \lipsum[1]
  \section*{Objetivos} 
  \lipsum[21]
  \section*{Metodologia} 
  \lipsum[3]
  \section*{Resultados e Discussão} 
  \lipsum[4]
  \section*{Considerações Finais} 
  \lipsum[5]

  \vspace{\baselineskip}  % Atenção! manter igual ao resumo
  \textbf{Palavras-chave:} Palavra-chave. Outra Palavra-chave composta. Bla.
\end{resumo}


\begin{abstract}
  Digital Twins (DTs) became a fundamental component of Industry 4.0, yet the current operation of DTs in industrial manufacturing 
  environments do not cover the coordination and management of multiple concurrent instances across the heterogeneous shop-floor. 
  Although the existing literature has architectural approaches to address the composition, orchestration, and infrastructure they 
  lack a coherent design that treats scalability and interoperability as first-class requirements. To cover this gap, this work 
  proposes an architecture for the deployment, orchestration and lifecycle management of multiple DT instances in the industrial 
  shop-floor environment. This architecture introduces a containerized unit called ``DTPod'' that encapsulates the state management 
  and data-driven inference logic to act as the digital counterpart of the physical asset. It also has a Twin Manager component 
  responsible for provisioning, monitoring, and decommissioning of the pod instances through a container orchestration platform. 
  The data collected from the physical assets is handled by a dedicated ingestion layer combining a message broker, schema registry, 
  and event router, with potential ontological modeling based via knowledge graphs. The semantic interoperability follows the concepts 
  of the Conceptual Interoperability Model (LCIM) and can be addressed through a standard data exchange pattern, such as OPC-UA. 
  Finally, a proof-of-concept instantiation is planned to evaluate the architecture under realistic conditions.
  \vspace{\baselineskip} 
  \textbf{Keywords:} Digital Twins. Orchestration. Lifecycle Management. Architecture.
\end{abstract}

\listoffigures*  % O * evita que apareça no sumário
\listoftables*
\listoflistings*  
\listofalgorithms*

\listadesiglas*[5em]

\begin{listadesimbolos}
  $\gets$   & Atribuição \\
  $\exists$   & Quantificação existencial \\
  $\rightarrow$   & Implicação \\
  $\wedge$   & E lógico \\
  $\vee$   & Ou lógico \\
  $\neg$   & Negação lógica \\
  $\mapsto$   & Mapeia para \\
  $\sqsubseteq$   & Subclasse (em ontologias) \\
  $\subseteq$   & Subconjunto: $\forall x\;.\; x \in A \rightarrow x \in B$ \\
  $\langle\ldots\rangle$ & Tupla \\
  $\forall$   & Quantificação universal \\
  mmmmm & Nenhum sentido, apenas estou aqui para demonstrar a largura máxima dessas colunas. Ao abrir o ambiente \texttt{listadesimbolos}, pode-se fornecer um argumento opcional indicando a largura da coluna da esquerda (o default é de 5em): \texttt{\textbackslash{}begin\{listadesimbolos\}[2cm] .... \textbackslash{}end\{listadesimbolos\}} \\
\end{listadesimbolos}

\tableofcontents*%

%%% Local Variables:
%%% mode: latex
%%% TeX-master: "main"
%%% End:

\textual%
% Assume-se que \textual já foi feito

\chapter{Exemplo}
\label{ch:exemplo}

\index{bobagem} Primeiro parágrafo da seção com uma frase sem sentido que só serve para ocasionar uma quebra e de demonstrar a configuração de indentação da primeira linha. Essa frase está aqui pois parágrafos de uma linha são feios.

Resultado do uso de siglas:
\begin{itemize}
\item Sigla que nunca expande: \API;
\item Sigla normal, expande no primeiro uso: \DHT, mas não no segundo: \DHT;
\item Siglas com plurais automaticos: \APIs e \DHTs;
\item Plural não-trivial: \SQs;
\item Forçando uma expansão (e no plural) \Glsfirstplural{DHT};
\item Usando uma sigla cujo comando é diferente da sigla: \WTC.
\end{itemize}

Resultado do glossário:
\begin{itemize}
\item Dois termos, \polling e \proxy;
\item Plural: \proxys.
\end{itemize}

Resultado de \mla|index|: primeiro um link normal \indexterm{tomate}, depois um capitalizado \indexTerm{tomate}.

Exemplo de autoref com \mla|amsthm|: \autoref{def:definition}, \autoref{def:defn}, \autoref{lemma:1}, \autoref{theorem:1}, \autoref{proof:1}.

\begin{definition}
  \label{def:definition}
  Exemplo de definição.
\end{definition}

\begin{defn}
  \label{def:defn}
  Exemplo de definição com ambiente defn.
\end{defn}

\begin{lemma}
  \label{lemma:1}
  Exemplo de Lema.
\end{lemma}

\begin{theorem}
  \label{theorem:1}
  Exemplo de teorema.
  \begin{proof}
    Exemplo de prova dentro do teorema \qed
  \end{proof}
\end{theorem}

\begin{theoremproof}
  \label{proof:1}
  Exemplo de prova \qed
\end{theoremproof}

(Sub)enumerações e citações (verificar se OK com o idioma):
\begin{enumerate}
\item \cite{turing1937}:
  \begin{enumerate}
  \item \citeonline{turing1937}:
    \begin{enumerate}
    \item \citeonline{dijkstra1968};
    \end{enumerate}
  \end{enumerate}
\item \cite{turing1937,dijkstra1968};
\item \citeonline{turing1937,dijkstra1968}.
\end{enumerate}


\begin{listing}[tb]
\caption{Meta informações do presente documento.}
\label{lst:meta}
\begin{minted}[highlightlines={1,4-5}]{latex}
\titulo{Template \LaTeX{} para testes e dissertações do LAPESD/UFSC}
\autor{André Roberto Alves de Oliveira}
\data{1 de agosto de 2019} % ou \today
\tese % ou \dissertacao
\titulode{Mestre em Ciência da Computação}
\orientador{Prof. Frank Augusto Siqueira, Dr.}
% \coorientador{Prof. Lars Thørväld, Dr.}

\membrobanca{Prof. Valerie Béranger, Dr.}{Universidade Federal de Santa Catarina}
\membrobanca{Prof. Mordecai Malignatus, Dr.}{Universidade Federal de Santa Catarina}
\membrobanca{Prof. Huifen Chan, Dr.}{Universidade Federal de Santa Catarina}
\coordenador{Prof. Charles Palmer, Dr.}
\end{minted}
\fonte{o autor.}
\end{listing}


Resultado de \mla|\autoref|s:
\begin{itemize}
\item \autoref{lst:meta};
\item \autoref{alg:algoritmo};
\item \autoref{fig:figura} tem subfiguras:
  \begin{itemize}
  \item \autoref{fig:svg}
  \item \autoref{fig:brasao}
  \end{itemize}
\item \autoref{tb:tabela};
\item \autoref{ch:exemplo};
\item \autoref{sec:frutas};
\item \autoref{sec:goiaba};
\item \autoref{sec:jabuticaba};
\item \autoref{sec:tomate}.
\end{itemize}


\begin{algorithm}
  \caption{Exemplo do ambiente \texttt{algorithimic}.}
  \label{alg:algoritmo}
  \begin{algorithmic}[1]
    \Procedure{Closure}{C, A}
      \State{$H \gets \emptyset$}\Comment{Direct cache}
      \For{$i \in [1, n]$}\Comment{Parallel, (dynamic,32) scheduling}
        \State{$H \gets H \cup \Call{DoImportantStuff}{i}$}
      \EndFor
    \EndProcedure
  \end{algorithmic}
  \fonte{o autor.}
\end{algorithm}

\begin{figure}[tb]
  \centering
  \caption{Exemplo de figura com duas subfiguras.}   
  \label{fig:figura}
  
  % Subfiguras são feitas usando as funcionalidades do memoir. Não
  % inclua outros pacotes, pois eles podem fazer o memoir dar ragequit
  % 
  % Há duas maneiras, a maneira limpinha (só no lapesd-thesis.cls) e a
  % maneira do memoir (aviso: \subtop não funciona direito).
  \subcaptionminipage[fig:svg]%
    {.49\linewidth}%
    {O Makefile compila SVGs em PDFs usando o inkscape}%
    % {\includegraphics[width=.2\linewidth]{alphachannel.pdf}}%
  \hfill% 
  % o comando acima expande para o equivalente disso:
  \begin{minipage}[t]{.49\linewidth}%
    \centering
    \subcaption{Brasão da UFSC.\label{fig:brasao}}
    \includegraphics[width=.2\linewidth]{\jobname-logo.pdf}
  \end{minipage}

  \fonte{o autor.}
\end{figure}

\begin{table}[tb]
  \centering
  \caption{Exemplo de tabela e símbolos}
  \label{tb:tabela}
  \begin{tabular}{lccp{5cm}}
    \toprule
    Esquerda & Coluna 1    & \rotatebox{90}{90 graus}  & Parágrafo com \mla|p{5cm}|   \\
    \midrule
    $r_1$    & \cmk        &  \xmk                     & \circledi    \\
    $r_2$    &     \multicolumn{2}{c}{merged cell}     & \circledii   \\
    $r_3$    & \circlediii & \circlediv                & \circledv    \\
    $r_4$    & \circledvi  & \circledvii               & \circledviii \\
    $r_5$    & \circledix  &  x                        & y           \\
    \bottomrule 
  \end{tabular}
  \fonte{o autor.}
\end{table}

\section{Introdução}
\subsection{Contextualização do Problema}
\subsection{Objetivos}  
\subsection{Proposta}
\subsection{Contribuições}
\subsection{Organização do Texto}

\section{Fundamentação Teórica}
\subsection{Gêmeo Digital}
\subsubsection{Conceitos}

The concept of Digital Twins (DT) originated from Michael Grieves' work on Product Lifecycle Management (PLM),
the original term was ``Conceptual Ideal for PLM", later he referred to as the ``Mirrored Spaces Model" and the conceptual model
was called ``Information Mirroring Model", the term currently used, Digital Twin, was coined at NASA \cite{Grieves2017}.\\
The first use of Digital Twin technology, considered by the academia, originated at NASA during the Apollo 13 incident.
The need to simulate the effects of the anomaly to understand the problem and find a way to bring the astronauts back safely,
through real-time testing to evaluate solutions, led the engineers in charge to modify the pilot training simulators to replicate
the current state of the lunar module \cite{Allen2021}.\\
Formally, a DT is defined as a virtual replica of a physical asset, process, or system, being a software composed of models 
and services, that intentionally, represent the original system throught its lifecycle \cite{Dalibor2022}.

\subsubsection{Classificação}

The concept of Digital Twins in the literature often suffers from a lack of distinction between its representations and different
levels of complexity. Although the concept of DTs was originally introduced by Grieves in 2002, the technical differentiation between
variations of virtual replicas was consolidated by \citeonline{Kritzinger2018}. As proposed by \citeonline{Kritzinger2018}, the differentiation
of the digital model occurs through the level of integration of the data flow between the physical and virtual worlds.
The representation of the models is categorized into three different levels: Digital Twin, Digital Model, and Digital Shadow.\\
The differentiation proposed by \citeonline{Kritzinger2018} establishes an information flow as the dividing line between the categories.
While manual flow characterizes the initial models, increasing automation and the way data is exchanged define the transition to the
terms ``shadow'' and ``digital''. This distinction is illustrated in the communication logic between physical and virtual assets,
as detailed below:
\begin{itemize}
    \item Digital Model
    \item Digital Shadow
    \item Digital Twin
\end{itemize}
%A diferença entre Digital Model, Digital Shadow e Digital Twin%
% Talvez seja interessante colocar a definição de 5 camadas do FULLER aqui tbm %

\subsection{Arquitetura de Referência e Ciclo de Vida}

%Falar desde o começo do ciclo sobre os sensores, depois dados e sobre os serviços e a interface%

\subsection{Tecnologias Habilitadoras}

%breve resumo sobre os enablers%

\subsection{Do Gêmeo Único ao Ecossistema}

%Citar os trabalhos que falam sobre a execução de multiplos DTs%

\subsection{Desafios de Execuções e Escalabilidade}

%falar dos problemas em outras pesquisas%
\section{Trabalhos Relacionados / Revisão da Literatura}

\section{Proposta}
\section{Metodologia}
\section{Cronograma}

\section{Frutas}
\label{sec:frutas}
\lipsum[4]

\subsection{Goiaba}
\label{sec:goiaba}
\lipsum[4]

\subsubsection{Jabuticaba}
\label{sec:jabuticaba}
\lipsum[4]

\subsubsubsection{Tomate}
\label{sec:tomate}

\begin{figure}[tb]
  \centering
  \caption{Segunda Figura.}
  \label{fig:segunda-fig}
  \includegraphics[width=.2\linewidth]{\jobname-logo.pdf}
  \fonte{o autor.}
\end{figure}

\xindex{tomate} \lipsum[4]


%%% Local Variables:
%%% mode: latex
%%% TeX-master: "main"
%%% End:

\postextual
\bibliography{main}
% Assume-se que \postextual já foi feito

\imprimirglossario

\apendices

\chapter{Systematic Literature Review Table}
\label{apx:systematic_literature_review_table}

\setlength{\LTpre}{10pt}

\footnotesize 

\begin{longtable}{
    >{\raggedright\arraybackslash\hspace{0pt}}p{1.8cm} % Autor
    >{\raggedright\arraybackslash\hspace{0pt}}p{1.5cm} % Setor (ajustado para quebrar texto)
    >{\raggedright\arraybackslash\hspace{0pt}}p{5.5cm} % Contribuição Principal
    >{\raggedright\arraybackslash\hspace{0pt}}p{2.7cm} % Tecnologias
    >{\raggedright\arraybackslash\hspace{0pt}}p{2.4cm} % Tipo de Impl.
}
    \caption{Summary of Related Works on Digital Twins}
    \label{tab:survey_dt} \\
    \toprule
    \textbf{Reference} & \textbf{Sector} & \textbf{Key contribution} & \textbf{Technology} & \textbf{Implementation type} \\
    \midrule
    \endfirsthead
    
    \multicolumn{5}{l}{\footnotesize\textit{Continuation of the Table \ref{tab:survey_dt}}} \\
    \toprule
    \textbf{Reference} & \textbf{Sector} & \textbf{Key contribution} & \textbf{Technology} & \textbf{Implementation type} \\
    \midrule
    \endhead
    
    \midrule
    \multicolumn{5}{r}{\footnotesize\textit{Continued on the next page.}} \\
    \endfoot
    
    \bottomrule
    \multicolumn{5}{l}{\footnotesize Source: The author (2026).} \\
    \endlastfoot
    
    % --- CONTEÚDO DA TABELA ---
    
    \citeonline{Juhlin2025} & Manufacturing & Proposed the CRADLE architecture, an interoperable DT architecture for cross-vendor integration & AAS, OPC-UA, AutomationML, Functional Mockup Interface, MQTT & Architectural Proposal + Framework + Prototype \\
    \citeonline{Pan2025} & Generic & Proposes a Digital Twin Network Architecture and a taxonomy system to define a DT fidelity, maturaty and complexity & - & Architectural Proposal + Survey \\
    \citeonline{Schroeder2021} & Generic & Methodology for implementing DTs and the concept of aggregated DTs, where one DT is composed of smaller DTs & AutomationML, Python, HTML, OpenGL, WatchDog, MATLAB, XML & Architectural Proposal + Proof of Concept \\
    \citeonline{Barbone2024} & Automotive/\hspace{0pt}Urban Mobility & Proposes the Digital Twin Continuum as a way to coordenate multiple DTs across edge-cloud environment & Docker, Kubernetes, White Label Digital Twin Library, Prometheus, Grafana, Helm & Prototype + Framework \\
    \citeonline{Alhazmi2025} & Generic/\hspace{0pt}Smart City & Layered architecture to providade data distribution in a MDTN contex & MQTT, JWTs & Architectural Proposal + Framework \\
    \citeonline{Mahoro2023} & Generic & Prposes Composite Digital Twin, a digital twin made of smaller ones to represent a complex system & Thing in The Future (Thing’In) & Architectural Proposal + Framework \\
    \citeonline{Merkle2019} & Energy/\hspace{0pt}Automotive Battery & Middleware framework to control and operate a fleet of DTs of eletric car batteries & AWS, MQTT, AWS S3, AWS DynamoDB, AWS Lambda, Python & Framework + Proof of Concept \\
    \citeonline{Qian2024} & Smart/\hspace{0pt}Generic & Proposes a distributed data layer to receive data from CPS and run multiple DTs & - & Architectural Proposal \\
    \citeonline{Picone2023} & Generic & Proposes an Edge Digital Twin, an architecture to run DTs in the edge & Java, MQTT, JSON, CoAP, Graphite, Eclipse Ditto (como comparador edge), Azure (como comparador cloud) & Architectural Proposal + Proof of Concept \\
    \citeonline{Tripathy2024} & Manufacturing & Proposal of a distributed solution architecture for condition monitoring and PLM based on DTs an architecture using the EAF & Eclipse Arrowhead Framework, EDMtruePLM & Architectural Proposal \\
    \citeonline{Wang2025} & Energy & Focus on emantic interoperability method based on ontology with multi fusion models for DTs & - & Framework \\
    \citeonline{Lai2024} & Smart Cities & Addresses the concept of multiple DTs by focusing on how individual primitive models can be aggregated and orchestrated to represent complex large-scale environments & Eclipse Ditto, Microsoft Azure lots of things from azure, MQTT, Kafka, InfluxDB, Grafana, MongoDB & Mapping Study + Architectural Proposal + Proof of Concept \\
    \citeonline{Maree2024} & Automotive & Presents an architecture for an IoT DTaaS plataform that facilitates the provisioning, management, and monitoring of heterogeneous IoT devices with DTAs and DTIs and TDTA & MQTT, gRPC & Architectural Proposal + Proof of Concept \\
    \citeonline{Martella2025} & Generic & Proposes a federation of DTs using the NGSI-LD standards & NGSI-LD, JSON-LD, Scorpio Context Broker, FIWARE Data Space Connector, Gaia-X Trust Framework & Architectural Proposal + Proof of Concept \\
    \citeonline{Monsalve2021} & Energy & A novel framework architecture for distributed digital twins in power systems using MAS & Python, Docker & Architectural Proposal + Framework \\
    \citeonline{Kuruppuarachchi2023} & Manufacturing & Proposes an architecture for Composite Digital Twin designed to enable and manage ecosystems in a secure manner & Docker, Hyperledger fabric, node.js, IPFS & Architectural Proposal \\
    \citeonline{Nie2023} & Manufacturing & Proposal for a cloud-edge orchestration architecture based on MAS and DT to optimize production control in distributed manufacturing environments & Python, Aliyun (Alibaba Cloud) & Architectural Proposal + Proof of Concept + Framework \\
    \citeonline{Nguyen2024} & Smart Cities & Proposes a 4 layer framework for DT orchestration & xData (NICT Japan), Eclipse Ditto, VMs, MQTT & Framework + Proof of Concept \\
    \citeonline{Xu2024} & Manufacturing & Proposal for a Cognitive Digital Twin framework designed for collaborative manufacturing with multiple robots (MRCMfg) & MuJoCo, ROS (Robot Operating System), URScript & Framework + Proof of Concept \\
    \citeonline{Ali2025} & Manufacturing & Proposes a blockchain based framework for cross-domain DTs & Hyperledge Fabric & Framework \\
    \citeonline{Cakir2025} & Energy & Proposes the Internet of Twins, a disaggregated DT deployed in a distributed cloud architecture and developed DT as a Plataform with knowledge graph based orchestration & gRPC, JSON, Python, Golang & Architectural Proposal + Proof of Concept \\
\end{longtable}

% Restaura o tamanho da fonte padrão para o restante do documento
\normalsize

\imprimirindice


%%% Local Variables:
%%% mode: latex
%%% TeX-master: "main"
%%% End:



\end{document}
